\subsection{All terminal pairs with minimum one at a unit root and a two-edge outlet}\label{proof:R2-C026}
For every terminal residual $(1,b)$, $b\ge2$, every inner length $L\ge1$, an internal root with exactly one residual unit side arm and an outer path of exactly two edges, the direction table below gives an actual output $U$ with $I_q(U)^3\le1$, where $I_q(V)=\E[q^{-1/3}|V|^{8/3}]$. Equal pairs remain at their original vertices.

\paragraph{R26.2 Original finite facts and the cubic certificate}

\paragraph{R26.2.1 The actual probability space and same-kernel moments}

Throughout this proof, \(ER^2\) means E[\(R^2\)]; the square of a mean is
written (\(ER)^2.\) All such expectations are on the same original mu.

Delete zero-mass atoms only. Then p(x), \(q(x)\ge \mu(x)>0\) because the diagonal
contribution has K(x, \(x)=0\) and all off-diagonal factors 1+/-K are nonnegative.
Thus inverse weights are legitimate without a uniform lower bound on p, q.

The C014 calculus gives, on the original L2(mu),
\[
\begin{gathered}
0\le d\le\tfrac13,\quad d^2\le \delta_2,\quad\\
\|K\|^2\le\theta=\tfrac{13}{32},\quad\\
\delta_{i+1}\le\theta \delta_i.
\end{gathered}
\]
Write \(k_i=K^i1,\allowbreak  r=k_2,\allowbreak  s=k_3,\allowbreak  t=k_4,\allowbreak  v=k_5.\)
Skew-adjointness gives
\[
Er=-d,\quad Et=\delta_2,\quad Eg=Es=Ev=0,
\]
\[
E(gr)=E(gt)=E(rv)=0,\quad E(gs)=-\delta_2,\quad E(gv)=\delta_3.
\]
The three further inequalities used below follow directly:
\[
\delta_4\ge \delta_2^2,\qquad d\delta_3\ge \delta_2^2,\qquad \delta_3\ge d\delta_2.
\]
The last follows from the second and \(\delta_2\ge d^2\) when \(d>0\); \(d=0\) is immediate.

For \(EV=0,\allowbreak \) subtract the row mean g(x) before Cauchy:
\[
\begin{gathered}
|KV(x)|^2\\
\le E_y(K(x,y)-g(x))^2\,EV^2\\
\le(1-g(x)^2)EV^2=p(x)q(x)EV^2.
\end{gathered}\tag{R26.A.1}
\]
For any original real R, a direct calculation gives the useful identity
\[
T^2R=(ER)q+E(gR)+K^2R,\qquad T^21=q+r=\bar f_2.     \tag{R26.A.2}
\]
In particular, a constant is propagated to \(\bar f_2,\allowbreak \). Also, for any
nonnegative \(R^2,\allowbreak \)
\[
T^2(R^2)\le 2E(R^2)q,                            \tag{R26.A.3}
\]
by first using \(T(R^2)\le 2E(R^2),\allowbreak \) then applying the positive T.

\paragraph{R26.2.2 Derivation of the entire cubic expansion}

Suppose \(h=aq+j\ge 0,\allowbreak  a>0,\allowbreak  j=KV,\allowbreak  EV=0,\allowbreak \) and write
\[
Ej=-\beta d,\quad C_1=E(gj)\ge0,\quad
N=EV^2,\quad J_2=Ej^2.
\]
The parameters a and beta need not agree. From (R26.A.1),
\[
E(j^4/q)\le N E(pj^2)\le2NJ_2.                  \tag{R26.A.4}
\]
Put
\[
W=2j-ag=K(2V-a),\qquad z=a+2N/a.
\]
The original row norm yields |\(W|\le \sqrt{a^2+4N}\le z,\allowbreak \) while
\[
EW=-2\beta d,\qquad EW^2=a^2d+4J_2-4aC_1.
\]
The pointwise polynomial (W+2 beta \(d)^2(W-z)\le 0\) therefore gives
\[
EW^3\le(z-4\beta d)EW^2-4z\beta^2d^2+8\beta^3d^3.
\]
All cubic terms in the exact \(h^4/q\) expansion combine as
\[
-a^4Eg^3+4a^3E(g^2j)-6a^2E(gj^2)+4aEj^3
=\frac a2\{EW^3+a^2E(g^2W)\}.
\]
Using \(E(g^2W)\le zd\) and (R26.A.4), we obtain
\[
\begin{gathered}
E(h^4/q)\le\mathcal A+(8a^2+6N-8a\beta d)J_2\\
-2a^2(4a+z-4\beta d)C_1,
\end{gathered}\tag{R26.A.5}
\]
where
\[
\begin{aligned}
\mathcal A={}&a^4+4a^3(a-\beta)d+2a^2Nd
-2a^3\beta d^2-2a^2\beta^2d^2\\
&-4N\beta^2d^2+4a\beta^3d^3\\
={}&a^4+2a^2d(N-\beta^2d)
+4a^3(a-\beta)d-4N\beta^2d^2\\
&-2a\beta d^2(a^2-2\beta^2d).
\end{aligned}                                                        \tag{R26.A.6}
\]
No fourth term or unfavorable cubic remainder is omitted, and the signed \(C_1\) term is kept.

Cauchy applied to \(Ej=-<g,\allowbreak  V>\) gives \(N\ge \beta^2\) d. Consequently
\[
z-4\beta d
\ge a\{1-2d+2d(\beta/a-1)^2\}>0.
\]
Under \(a\le \beta\le 6a/5\) and \(N\le 4a\) beta d/3, every remainder discarded from
(R26.A.5)-(R26.A.6) is nonpositive: in particular
\[
a^2-2\beta^2d\ge a^2/25,\qquad
6N-8a\beta d\le0.
\]
This gives
\[
E(h^4/q)\le a^4+8a^2J_2+2a^2d(N-\beta^2d).       \tag{R26.A.7}
\]
If \(a\le \beta\) is not available, retain \(4a^3(a-\beta)_+d\); this is precisely
the positive-part variant of the displayed certificate.

For the broader variant, assume only \(0<a\le 1,\allowbreak  0<\beta\le 1,\allowbreak  C_1\ge 0\) and
\[
(a-\beta)_+\le w,\qquad N\le nd,\qquad J_2\le jd.
\]
Dropping only the explicit negative terms of (R26.A.6), and retaining the
positive cubic term as \(4ad^3,\allowbreak \) gives
\[
E(h^4/q)\le
a^4+(4wa^3+8ja^2)d+(2na^2+6nj)d^2+4ad^3.       \tag{R26.A.8}
\]
This is the broader variant above.

The following coefficient formula provides an independent check on the
later displayed polynomials. Substitution \(a=1-\rho\) d in (R26.A.8) has
coefficients, after the constant 1,
\[
\begin{array}{c|l}
d&-4\rho+4w+8j\\
d^2&6\rho^2-12w\rho-16j\rho+2n+6nj\\
d^3&-4\rho^3+12w\rho^2+8j\rho^2-4n\rho+4\\
d^4&\rho^4-4w\rho^3+2n\rho^2-4\rho .
\end{array}                                                        \tag{R26.A.9}
\]
The comparisons in the sequel substitute into the monotone expression
(R26.A.8) before using the resulting signed polynomial.

\paragraph{R26.3 New row \(b=2,\allowbreak  L=1\)}\label{part:R26.3}

Set \(\gamma=d-\delta_2/2,\allowbreak  c=1-\gamma/2\) and \(\beta=c^2.\) Since \(p+\bar f_2=2+r,\allowbreak \)
\[
p\bar f_2\le(1+r/2)^2,\qquad
A(p\bar f_2)\le p-\gamma+s.
\]
Apply Corollary~\ref{cor:unit-two-step:first-square} to the inner
operator $A=J+K$ (inner kernel $-K$ in the common notation). It gives
\[
qA(p\bar f_2)\le(c+s/2)^2.
\]
Using \(T^2s=-\delta_2+v,\allowbreak \) (R26.A.2)-(R26.A.3) give
\[
0<U_*\le h=aq+j,
\]
\[
\begin{gathered}
a=\beta-c\delta_2+\delta_3/2,\quad\\
V=\beta g+c(t-\delta_2),\quad\\
j=KV=\beta r+c(v-\delta_2g).
\end{gathered}\tag{R26.A.10}
\]
The exact covariance data are
\[
Ej=-\beta d,\quad C_1=c(\delta_3-d\delta_2)\ge0,
\]
\[
\begin{gathered}
N=\beta^2d+c^2(\delta_4-\delta_2^2),\\
\quad\\
J_2=\beta^2\delta_2+c^2(\delta_5-2\delta_2\delta_3+\delta_2^2d)\\
\le\beta^2\delta_2+c^2\delta_5.
\end{gathered}
\]
Both omitted cross terms vanish by skew parity, rather than by an
independence assumption.

The certificate hypotheses are paid as follows:
\[
\begin{gathered}
\tfrac56\le c\le1,\quad\tfrac{25}{36}\le\beta\le1,\quad\\
\tfrac a\beta\ge\tfrac{67}{80}>\tfrac56,\quad a\le\beta,
\end{gathered}
\]
\[
\begin{gathered}
\frac{N}{a\beta d}\\
\le\frac65\left(1+\frac{36}{25}\theta^3\right)\\
<\frac{831}{625}<\frac43.
\end{gathered}
\]
The divisions are used only for \(d>0.\) Also \(N\ge \beta^2d\) is immediate from
its exact expression.

Applying (R26.A.7), \(\delta_5\le \theta^3\delta_2\) and \(\delta_4\le \theta^3d\) yields
\[
\begin{gathered}
E(h^4/q)\le\\
a^4+a^2\left(\frac{559}{160}d+\frac3{20}d^2\right).
\end{gathered}\tag{R26.A.11}
\]
The exact scalar is
\[
a=1-d-\frac12\delta_2+\frac14d^2+\frac14d\delta_2-\frac3{16}\delta_2^2+\frac12\delta_3.
\]
Using \(d^2\le \delta_2\) and \(d\le 1/3\) gives
\[
a\le1-d+\frac7{192}\delta_2
\le1-\frac{6053}{6144}d
\le1-\frac{63}{64}d.
\]

For \(0\le z<1,\allowbreak \)
\[
(1-z)^{-2}\ge1+2z+3z^2.
\]
With \(\rho=63/64,\allowbreak  C=559/160,\allowbreak  D=3/20,\allowbreak \) this proves
\[
\begin{gathered}
1-(1-\rho d)^2\{(1-\rho d)^2+Cd+Dd^2\}\\
\ge(1-\rho d)^2\{(4\rho-C)d+(2\rho^2-D)d^2\}.
\end{gathered}
\]
Here \(4\rho-C=71/160,\allowbreak  2\rho^2-D>0\) and 1-rho \(d\ge 43/64.\)
Thus the loss is at least
\[
\frac{131279}{655360}d>\frac15d.
\]
It follows that \(E(U_*^4/q)\le 1-d/5\) on the entire legal host domain.


\paragraph{R26.4 The new finite inner lengths \(L=2,\allowbreak \) 3, 4}

\paragraph{R26.4.1 Two same-measure difference-moment facts}

Since \(E(r+d)=0,\allowbreak \)
\[
\delta_3-2d\delta_2+d^3=\|K(r+d)\|^2\le\theta(\delta_2-d^2).
\]
Subtracting this bound from the desired upper expression leaves
(\(1-\theta-d)(\delta_2-d^2)\ge 0.\) Hence
\[
\begin{gathered}
\delta_3\le(1+d)\delta_2-d^2,\qquad\\
\delta_2-\delta_3\ge d(d-\delta_2).
\end{gathered}\tag{R26.A.12}
\]
Let sigma be the spectral measure of g for -\(K^2.\) Cauchy for the
nonnegative measure (1-x)d sigma gives
\[
(\delta_2-\delta_3)^2\le(d-\delta_2)(\delta_3-\delta_4).
\]
For \(d>0,\allowbreak \) (R26.A.12) then gives
\[
\delta_3-\delta_4\ge d(\delta_2-\delta_3).                            \tag{R26.A.13}
\]
These are actual moment inequalities. They are not the false assertion
that every original spectral point satisfies \(x\ge d.\)

\paragraph{R26.4.2 \(L=2\)}

Put
\[
B=1-d+\delta_2/2,\quad\tau=d-\delta_2,\quad R=r+t,\quad
V_R=ER^2=\delta_2-2\delta_3+\delta_4.
\]
Positive propagation from the preceding row yields
\[
A^2(p\bar f_2)\le Bp+R,\qquad
qA^2(p\bar f_2)\le B(1+R/(2B))^2.
\]
Using \(T^2r=-dq+t\) and \(T^2t=\delta_2q+k_6\) gives the actual upper bound
\[
\begin{gathered}
U_*\le aq+j,\quad a=B-\tau+V_R/(2B),\\
\quad V=Bg+s+v,\quad j=KV=Br+t+k_6.
\end{gathered}\tag{R26.A.14}
\]
For \(d>0\) write \(x=\delta_2/d,\allowbreak  y=\delta_3/d.\) The odd filter is
\(phi(z)=B-z+z^2,\allowbreak \) and \(\beta=B-x+y.\) On \(0\le z\le \theta,\allowbreak  0<phi(z)\le B\le 1.\) Therefore
\[
\begin{gathered}
Ej=-\beta d,\quad C_1=0,\quad\\
N=\int\phi^2\,d\sigma,\quad J_2=\int z\phi^2\,d\sigma,
\end{gathered}
\]
\[
\beta^2d\le N\le d,\qquad J_2\le\theta d.
\]
Equation (R26.A.12) gives \(a\ge \beta>0.\) Also \(V_R\le \delta_2\) and
\(2Btau>\delta_2\) show \(0\le a-\beta\le x\le \theta.\)

To verify the nontrivial certificate condition, set \(\kappa=19/32\) and
\(m=B-\theta+\theta^2\ge B/2.\) Since \(B-phi(z)\ge \kappa\) z,
\[
N\le B\beta d-\kappa m \delta_2
\le\beta d(B-\kappa x/2)
\le\beta d(B-\kappa d/2).
\]
The exact sufficient comparison is
\[
4a-3B+\tfrac32\kappa d
\ge1-\frac{263}{64}d+\frac92d^2\ge\frac{25}{192}>0.
\]
It follows that \(N\le 4a\) beta d/3; \(\beta/a\le 1\) is already known.

For the scalar loss, log-convexity and \(\delta_4\le \theta \delta_3\) give
\[
V_R\le dx\{1-(2-\theta)x\},\qquad B\ge(4+x)/6,
\]
so
\[
a\le1-d\left[2-\frac32x-
\frac{3x\{1-(2-\theta)x\}}{4+x}\right]\le1-\frac54d.
\]
The final inequality is equivalent to \(264x-105x^2\le 96.\) Its left side is
increasing on [0, theta], and the endpoint margin is \(6225/1024.\)

The retained-remainder form of (R26.A.7) now gives
\[
E(h^4/q)\le a^4+a^2\left(\frac{39}{8}d+2d^2\right).
\]
The same reciprocal-square comparison as in \mbox{\hyperref[part:R26.3]{part R26.3}}, now with
\(\rho=5/4,\allowbreak  C=39/8,\allowbreak  D=2,\allowbreak \) supplies loss
\[
(7/12)^2(1/8)d=\frac{49}{1152}d>\frac1{24}d.
\]
Thus \(E(U_*^4/q)\le 1-d/24.\)

\paragraph{R26.4.3 \(L=3,\allowbreak \) including the propagated constant}\label{part:R26.4.3}

Let \(\omega_2=\delta_2-\delta_3\) and retain B, tau. This time
\[
A^3(p\bar f_2)\le Bp+R-\tau,\qquad R=Br+s+v.
\]
The root upper envelope is \(B(1+(R-\tau)/(2B))^2.\) With
\[
V_R=E(R-\tau)^2
=B^2\delta_2+\delta_3-2\delta_4+\delta_5+2\tau Bd+\tau^2,
\]
the exact propagated upper bound is
\[
\begin{aligned}
a&=B(1-d)-\tau-\omega_2+V_R/(2B),\\
j&=(B-\tau)r+Bt+v+k_7-\omega_2g,\\
V&=(B-\tau)g+Bs+(t+k_6-\omega_2),\qquad j=KV.
\end{aligned}                                                        \tag{R26.A.15}
\]
The coefficient of r is B-tau, as it must be because the constant -tau
propagates to -tau \(\bar f_2.\) Replacing it by -tau q would change this row.

The odd filter \(phi(z)=B(1-z)-\tau\) is between \(1/16\) and 1. Its mean
coefficient is \(\beta=(B-d)(1-x)>0,\allowbreak \) and
\[
C_1=\delta_3-\delta_4-d(\delta_2-\delta_3)\ge0
\]
by (R26.A.13). The odd and centered even components of V are orthogonal.
Their squared norms give
\[
N\le d+\delta_4\le\frac{43}{40}d,\qquad
J_2\le\frac{559}{1280}d.
\]
Since \(h\ge U_*>0\) and \(Eh=a-\beta\) d, \(a>0.\)

The exact scalar expansion is
\[
a=1-3d+2d^2+(1-2d+\delta_2/4)\delta_2+\delta_3
+\frac{\delta_3-2\delta_4+\delta_5+\tau^2}{2B}.
\]
Use \(\delta_3-2\delta_4+\delta_5\le 4\delta_2/27\) and \(B\ge 2/3.\) After collecting the d terms before
substituting \(d\le 1/3,\allowbreak \) one obtains
\[
\begin{gathered}
\frac{a-1+3d}{d}\\
\le\frac{11}{12}+(\theta-1/18)x+\frac13x^2\\
\le\frac{325}{288}<\frac54.
\end{gathered}
\]
Thus \(a\le 1-7d/4.\) Also
\[
a-\beta\le Bx+d(1/4-3x/2).
\]
For \(x\ge 1/6\) this is at most x; for \(x<1/6,\allowbreak  d\le x\) makes it at most 5x/4.
Hence (\(a-\beta)_+\le \theta.\)

Apply (R26.A.8) with
\[
\rho=\tfrac74,\quad w=\theta,\quad n=\tfrac{43}{40},
\quad j=\tfrac{559}{1280}.
\]
Formula (R26.A.9) gives exactly
\[
\begin{gathered}
E(h^4/q)\le\\
1-\frac{301}{160}d+\frac{66111}{25600}d^2\\
+\frac{1707}{2560}d^3+\frac{651}{2560}d^4.
\end{gathered}
\]
The three positive coefficients are respectively less than 3, 1, 1.
Consequently the loss is at least
\[
\begin{gathered}
\left(\frac{301}{160}-1-\frac19-\frac1{27}\right)d\\
=\frac{3167}{4320}d>\frac7{10}d.
\end{gathered}
\]
This proves the \(L=3\) row, retaining every cubic and quartic term.

\paragraph{R26.4.4 \(L=4\)}

Set \(C=B-\tau=1-2d+3\delta_2/2\ge 1/2\) and
\[
R=B(r+s)+t+k_6.
\]
The true propagated envelope is
\[
A^4(p\bar f_2)\le Cp+R-Bd,\qquad
qA^4(p\bar f_2)\le C(1+(R-Bd)/(2C))^2.
\]
With \(V_R=E(R-Bd)^2,\allowbreak \) (R26.A.2) gives
\[
\begin{aligned}
a&=C-2Bd+(1-B)\delta_2-\delta_3+V_R/(2C),\\
j&=(C-Bd)r+Bt+k_6+k_8+B(v-\delta_2g),\\
V&=(C-Bd)g+Bs+v+k_7+B(t-\delta_2).
\end{aligned}                                                        \tag{R26.A.16}
\]
Again, the constant -Bd contributes -Bd \(\bar f_2,\allowbreak \) including its r term.

The odd filter
\[
\phi(z)=C-Bd-Bz+z^2-z^3
\]
decreases on [0, theta], because phi'\ensuremath{\le}-\(B+1/3<0.\) Its endpoint has the
strict original-parameter lower bound
\[
\begin{gathered}
\phi(\theta)\ge\\
(1-\theta)-(3-\theta)d+2d^2+\theta^2(1-\theta)\\
\ge\frac{14563}{294912}>0.
\end{gathered}
\]
It is also at most 1. The centered even component gives
\[
\begin{gathered}
0<\beta\le1,\quad\\
C_1=B(\delta_3-d\delta_2)\ge0,\quad\\
N\le\frac{43}{40}d,\quad J_2\le\frac{559}{1280}d.
\end{gathered}
\]

The even filter in R is \(B-z+z^2,\allowbreak \) between 0 and B, so
\[
ER^2\le B^2(\delta_2+\delta_3),\quad ER=-Bd+\delta_2-\delta_3,
\]
\[
V_R\le B^2(\delta_2+\delta_3)+3B^2d^2-2Bd(\delta_2-\delta_3).
\]
Since \(B^2-C\le d^2/2\) and \(C\ge 1/2,\allowbreak  B^2/C\le 1+d^2\le 10/9.\) Substitution gives
\[
\begin{gathered}
a\le1-4d+\frac{11}{3}d^2+\frac{37}{18}\delta_2\\
-\frac49\delta_3-\frac12\delta_2^2\\
\le1-\frac{373}{192}d\le1-\frac74d.
\end{gathered}
\]
Writing \(\beta=C-Bd-Bx+y-z\) with \(y=\delta_3/d,\allowbreak  z=\delta_4/d\) and retaining \(z\le \theta\) y,
\[
a-\beta\le x+d(5/4-3x/2)
\le\frac94x-\frac32x^2
\le\frac{1365}{2048}<\frac23.
\]

Apply (R26.A.8) with the previous n, j, rho and \(w=2/3.\) The exact polynomial is
\[
\begin{gathered}
1-\frac{403}{480}d-\frac{73889}{25600}d^2\\
+\frac{26207}{2560}d^3-\frac{40922}{7680}d^4.
\end{gathered}
\]
The last three coefficients are bounded above by -2, 11, 0.
The convex function -\(2d+11d^2\) has maximum \(5/9\) on [0, \(1/3\)]. Therefore
\[
E(h^4/q)\le1-\frac{409}{1440}d\le1-\frac14d.
\]
This completes all finite inner lengths \(L=1,\allowbreak \) 2, 3, 4 without a change of
host-dependent direction.

\paragraph{R26.5 New row \(b=3,\allowbreak  L=1\)}

The genuine three-edge inward arm is
\[
\bar f_3=q+r-s-d.
\]
Both the scalar -d here and both propagated \(\delta_2\) terms below are retained.
Put \(m=1-d/2,\allowbreak  S=s-t,\allowbreak  D=\delta_2+\delta_3.\) From
\(p+\bar f_3=2m+r-s\) and positivity,
\[
p\bar f_3\le(m+(r-s)/2)^2,\qquad
A(p\bar f_3)\le m^2p-md+mS+D/2.
\]
The same first-square corollary, with inner kernel $-K$, gives
\[
qA(p\bar f_3)\le(c+S/2)^2,\qquad
c=1-d+\frac{\delta_2+\delta_3}{4m}.
\]
The correct two-step identity is
\[
T^2S=-\delta_2-\delta_2q+v-k_6.
\]
Consequently
\[
\begin{gathered}
a=c^2-2c\delta_2+(\delta_3+\delta_4)/2,\quad\\
V=c^2g-cv+c(t-\delta_2),\quad\\
j=c^2r-ck_6+c(v-\delta_2g).
\end{gathered}\tag{R26.A.17}
\]
The odd filter is \(c^2-cz^2.\) The bounds
\[
\begin{gathered}
1-d\le c\le1-d+\frac{27}{64}\delta_2\le1,\qquad\\
c\ge1-d+(d^2+d^3)/4\ge\frac{19}{27}>\frac7{10}
\end{gathered}
\]
make it positive. Thus
\[
\begin{gathered}
\beta=c^2-c\delta_3/d>0,\quad\\
C_1=c(\delta_3-d\delta_2)\ge0,\quad\\
N\le\frac{43}{40}d,\quad J_2\le\frac{559}{1280}d.
\end{gathered}
\]
Also \(a\le c^2,\allowbreak \) (\(a-\beta)_+\le \theta^2\le \theta,\allowbreak \) and positivity of h gives \(a>0.\)

Using \(\lambda=1-c,\allowbreak \) the exact available scalar losses yield
\[
a\le1-2d+d^2-\frac{2771}{10240}\delta_2
\le1-\frac{53971}{30720}d\le1-\frac74d.
\]
The parameters of (R26.A.8) are therefore exactly those of \mbox{\hyperref[part:R26.4.3]{part R26.4.3}}.
Its already checked polynomial proves
\[
E(U_*^4/q)\le1-\frac7{10}d.
\]
This is a separate actual-message calculation for terminal (1, 3), not a
reuse of the terminal (1, 2), \(L=3\) message.

\paragraph{R26.6 Every original inner length \(L\ge 5\) with terminal (1, 2)}

\paragraph{R26.6.1 Finite convolution and the actual two losses}

Define
\[
\begin{gathered}
\zeta_1=\zeta_2=0,\quad\\
\zeta_j=\sum_{i=1}^{\lfloor(j-1)/2\rfloor}(-1)^{i-1}\delta_i,\quad\\
S_L=\sum_{j=0}^Lk_j.
\end{gathered}
\]
The identity \(S_j=A S_{j-1}\)+\(\zeta_j\) follows by taking the actual mean of
\(S_{j-1}\). Iteration gives
\[
f_L=S_L-\sum_{j=3}^L\zeta_jf_{L-j},\qquad
\tau\le\zeta_j\le d.
\]
The latter bounds are finite alternating-sum bounds, with \(\tau=d-\delta_2.\)
The finite Duhamel formula applied to r gives
\[
A^Lr=k_{L+2}-df_{L-1}+\delta_2f_{L-3}-\delta_3f_{L-5}+\cdots,
\]
with nonnegative indices only. Combining this with \(p\bar f_2\le (1+r/2)^2\)
and \(A(r^2)\le 2\delta_2\) gives
\[
\begin{gathered}
A^L(p\bar f_2)\le S_L+k_{L+2}-\gamma f_{L-1}\\
-\sum_{j=3}^L\omega_jf_{L-j},
\end{gathered}\tag{R26.A.18}
\]
where \(\omega_j=\zeta_j\) for even j and \(\omega_j= \zeta_{j+2}\) for odd j.
Every \(\omega_j\) is at least \(\tau\). This verifies the truncation-end signs.

For \(L\ge 5\) the \(j=L\) and \(j=L-1\) terms exist. Retain their tau multiples of
\(f_0\) and \(f_1\) and discard only other nonnegative subtractions. With
\[
B=1-\tau,\qquad R=\sum_{j=2}^Lk_j+k_{L+2},
\]
the resulting genuine nonnegative upper bound is
\[
\begin{gathered}
A^L(p\bar f_2)\le Bp+R-\tau,\\
\qquad H_*\le B(1+(R-\tau)/(2B))^2.
\end{gathered}\tag{R26.A.19}
\]
These are two losses already present in the original path expansion;
they are not additional tree arms.

\paragraph{R26.6.2 Every finite mask, rather than a length cutoff}

Let \(m=L-1\ge 4\) and \(P_m(z)=G_m(z)+ z^{m+1}\), where
\(G_m(z)=1+...+ z^{m-1}\). Then \(R=KP_m(K)g.\) For \(z=i\) sqrt(x),
\[
|P_m(z)|^2=
\frac{1+x^m(1-x+x^2)-2\Re[z^m(1-x-z)]}{1+x}.     \tag{R26.A.20}
\]
This identity is obtained by putting the whole finite numerator over 1-z;
it includes the final extra monomial.

The four residue classes of m give a complete proof:


\begingroup\small
\setlength{\tabcolsep}{3pt}
\setlength{\LTleft}{0pt}\setlength{\LTright}{\fill}
\begin{longtable}{@{}p{0.1200\linewidth}p{0.8000\linewidth}@{}}
\caption{All terminal pairs with minimum one at a unit root and a two-edge outlet: parameter values.}\label{tab:parameters:r2-c026:1}\\
\toprule
Entry & Numerator bound in (R26.A.20) \\
\midrule\endfirsthead
\toprule
Entry & Numerator bound in (R26.A.20) \\
\midrule\endhead
\bottomrule\endfoot
0 & At most 1; the subtraction is \(2x^{m/2}(1-x).\). \\
\addlinespace[2pt]
1 & At most 1; the subtraction is \(2x^{(m+1)/2}.\). \\
\addlinespace[2pt]
2 & At most \(1+2x^3+x^6,\allowbreak \) because \(m\ge 6.\). \\
\addlinespace[2pt]
3 & At most \(1+2x^4+x^7,\allowbreak \) because \(m\ge 7.\). \\
\end{longtable}
\endgroup


For the last two rows, \(2\theta^2+\theta^5<1\) proves |\(P_m|\le 1.\)
For the second weighted bound the only larger required expression is
\[
2x+6x^4+3x^7
\le\frac{13}{16}+\frac16+\frac{13}{2304}
=\frac{2269}{2304}<1.
\]
The other bad-residue expression is smaller. The first two rows use
\(x/(1+x)\le 13/45<1/3.\) Thus every finite \(m\ge 4\) satisfies
\[
\begin{gathered}
|P_m(i\sqrt x)|^2\le1,\qquad\\
x|P_m(i\sqrt x)|^2\le\tfrac13.
\end{gathered}\tag{R26.A.21}
\]


For \(A_k(x)=\sum_{i=0}^{k-1}(-x)^i\),
\[
1-x\le A_k(x)\le1\quad(k\ge1),
\]
and the exact real parts are
\[
\begin{array}{c|cc}
m&\Re P_m&\Re(i\sqrt xP_m)\\ \hline
2h&A_h&-xA_{h+1}\\
2h+1&A_{h+2}&-xA_h .
\end{array}
\]
The relevant indices are at least two. Therefore, on the same sigma,
\[
\begin{gathered}
-d\le ER\le-\tau,\quad E(gR)\le-(\delta_2-\delta_3),\\
\quad ER^2\le d/3.
\end{gathered}\tag{R26.A.22}
\]

\paragraph{R26.6.3 The covariance is on the spectral measure of 1}

It is necessary to retain the correct spectral measure in the sign
argument. For every finite k,
\[
A_k'(x)\ge-\frac1{(1-x)^2},
\]
so, on [0, theta],
\[
\begin{gathered}
\frac{d}{dx}(x^2A_k(x))\\
\ge x\left(2\kappa-\frac{\theta}{\kappa^2}\right)\ge0,\\
\quad 2\kappa^3-\theta=\frac{406}{32768}>0.
\end{gathered}\tag{R26.A.23}
\]
Let nu be the original probability spectral measure of 1 for -\(K^2.\)
Then d \(\sigma=x\) d nu and integral x d \(\nu=d.\) For increasing F,
\[
\begin{gathered}
\int(x-d)F(x)\,d\nu\\
=\tfrac12\iint(x-y)(F(x)-F(y))\,d\nu(x)d\nu(y)\ge0.
\end{gathered}\tag{R26.A.24}
\]
The legitimate choice is \(F=x^2A_k.\) This proves an integrated covariance
sign even when nu has spectral mass below d.

\paragraph{R26.6.4 Actual \(T^2\) propagation and its scalar payment}

Put \(\alpha=B-\tau=1-2\tau\) and \(V_R=E(R-\tau)^2.\) Since \(\tau\le 2/9,\allowbreak  \alpha\ge 5/9.\) Expanding (R26.A.19) with (R26.A.2) gives
\[
\begin{aligned}
a&=\alpha+ER+E(gR)+V_R/(2B),\\
j&=\alpha r+K^2R+E(gR)g,\\
V&=\alpha g+KR+E(gR),\qquad EV=0.
\end{aligned}                                                        \tag{R26.A.25}
\]
With \(\beta=\alpha- d^{-1}\) integral x \(Re(P_m)d\) sigma,
\[
Ej=-\beta d,\quad 0<5/9-\theta\le\beta\le1.
\]
The exact covariance is
\[
C_1=E[(s+dg)R]
=\int x^2(x-d)A_\ell(x)\,d\nu(x)\ge0
\]
by (R26.A.23)-(R26.A.24).

The vector V is the centered version of
(\(\alpha+K^2P_m(K))g.\) By \(Re(P_m)\ge \kappa>1/2\) and (R26.A.21),
\[
\begin{aligned}
N&\le\int|\alpha-xP_m|^2\,d\sigma\\
&\le\int(\alpha^2-\alpha x+x^2)\,d\sigma
\le\alpha^2d\le d,\qquad J_2\le\theta d,
\end{aligned}
\]
because \(x\le \theta<\alpha.\) Positivity of h again gives \(a>0.\)

Combining the same ER terms before discarding signs is essential:
\[
a=\alpha+\frac{\alpha}{B}ER+E(gR)
+\frac{ER^2+\tau^2}{2B}.
\]
Equation (R26.A.22) gives
\[
\begin{gathered}
a\le1-3\tau-(\delta_2-\delta_3)\\
+\frac{d/6+(3/2)\tau^2}{B}.
\end{gathered}\tag{R26.A.26}
\]
In particular,
\[
a-\alpha\le
\frac{-\tau+(5/2)\tau^2+d/6}{B}
\le-\frac{7d}{72B}\le0,
\]
so (\(a-\beta)_+\le \theta.\)

For \(x=\delta_2/d,\allowbreak  y=\delta_3/d,\allowbreak \) (R26.A.26) also gives
\[
a\le1-d\left[
3-2x-y-\frac{1/6+(3/2)d(1-x)^2}{1-d(1-x)}
\right]\le1-\frac32d.
\]
To verify the last inequality over the entire domain, use \(y\le \theta\) x,
the monotonicity of the displayed fraction in d, and \(d\le 1/3.\)
After multiplication by 2+x, the remaining expression is
\[
1-\frac5{16}x-\frac{125}{32}x^2
\ge\frac{7483}{32768}>0.
\]

In (R26.A.8), take \(\rho=3/2,\allowbreak  w=\theta,\allowbreak  n=1,\allowbreak  j=\theta.\) The full polynomial is
\[
E(h^4/q)\le
1-\frac98d+\frac78d^2+\frac{89}{32}d^3-\frac{123}{64}d^4.
\]
Discard the negative quartic term and use \(d^2\le d/3,\allowbreak  d^3\le d/9.\) The loss is
at least \(151d/288>d/2.\) Thus, for every original integer \(L\ge 5,\allowbreak \)
\[
E(U_*^4/q)\le1-\frac12d.
\]

\paragraph{R26.7 Every original terminal length \(b\ge 4\) with \(L=1\)}

\paragraph{R26.7.1 The finite arm and its genuine envelope}

From the finite convolution for T, retaining only its last nonnegative
loss,
\[
\bar f_b\le\sum_{j=0}^b(-K)^j1-\zeta_b.
\]
Let
\[
n=b-1\ge3,\quad m=1-\zeta_b/2,\quad
R=KG_n(-K)g=\sum_{j=2}^b(-K)^j1,
\]
\[
t_0=-ER=\zeta_{b+1},\quad V_R=ER^2,\quad S=KR,\quad
c=m-t_0/2+V_R/(4m).
\]
The sum \(p+\bar f_b\) is bounded by 2m+R, which is nonnegative because it
dominates that nonnegative sum. Thus
\[
p\bar f_b\le(m+R/2)^2,\quad
A(p\bar f_b)\le m^2p-mt_0+mS+V_R/2,
\]
\[
0<H_*\le(c+S/2)^2.                              \tag{R26.A.27}
\]

For every \(n\ge 3\) the finite geometric sum has |\(G_n(i \sqrt{x})|\le 1\) on
[0, theta]. Odd n gives numerator \(1+x^n\); n divisible by four gives
(\(1-x^{n/2})^2\); the remaining \(n\ge 6\) satisfies
\(2x^{n/2-1}+x^{n-1}\le 2\theta^2+\theta^5<1.\) Therefore
\[
ER^2\le \delta_2,\qquad ES^2\le \delta_3.                 \tag{R26.A.28}
\]
Writing
\[
\begin{gathered}
G_n(-i\sqrt x)=B_0(x)-i\sqrt xA_0(x),\\
\quad B_0=A_{\lceil n/2\rceil},\quad A_0=A_{\lfloor n/2\rfloor},
\end{gathered}
\]
both \(A_0,\allowbreak  B_0\) lie in [1-x, 1]. The exact means are
\[
e=-ES=E(gR)=\int xA_0\,d\sigma,\quad
\ell=-E(gS)=\int xB_0\,d\sigma,
\]
\[
\delta_2-\delta_3\le e,\ell\le \delta_2.                       \tag{R26.A.29}
\]
In this row E(gR) is positive, unlike the long-inner-path row. The
calculation retains this sign change.

The same finite arm supplies
\[
1-d+\delta_2/2\le c\le1-d+\frac{13}{10}\delta_2
\le1-\frac{151}{320}d.
\]
For the lower bound, when \(b=4\) the two zeta terms sum to \(2d-\delta_2.\)
For \(b\ge 5\) their sum is at most \(2d-2\delta_2+2\delta_3\le 2d-\delta_2.\)
Thus \(c\ge 13/18\); in particular \(c>\theta^2.\)

\paragraph{R26.7.2 The necessary variance improvement survives}\label{part:R26.7.2}

Use \(T^2S=-e q-\ell+K^2S\) in (R26.A.27). The actual nonnegative upper bound is
\[
\begin{aligned}
a&=c^2-c(e+\ell)+ES^2/2,\\
j&=c^2r+cK^2S-c\ell g,\\
V&=c^2g+cKS-c\ell,\qquad EV=0,\quad j=KV.
\end{aligned}                                                        \tag{R26.A.30}
\]
Its odd Krylov component is (\(c^2-cx^2A_0(x))g.\) Its centered even component
is \(c(K^3B_0(-K^2)g-\ell).\) They are orthogonal. Consequently, with
\(\xi= d^{-1}\) integral \(x^2A_0\) d sigma,
\[
Ej=-\beta d,\quad\beta=c^2-c\xi,\quad0<\beta\le1.
\]
The same original spectral covariance proof (R26.A.23)-(R26.A.24) gives
\[
C_1=c\int x^2(x-d)B_0(x)\,d\nu(x)\ge0.
\]

The stronger variance bound is not optional. Put \(z=x^2A_0(x),\allowbreak \) so
\(0\le z\le \theta^2<c.\) Then
\[
(c^2-cz)^2\le c^4-c^3z.
\]
Using \(A_0\ge \kappa\) and the centered even variance at most \(c^2\delta_4,\allowbreak \)
\[
N\le c^4d-c^3\kappa \delta_3+c^2\delta_4
\le c^4d+c^2(\theta-c\kappa)\delta_3\le c^4d,
\]
since c \(\kappa\ge 247/576>\theta=234/576.\)
For \(\lambda_0=151/320,\allowbreak \) the complete fourth expansion, with \(d\le 1/3,\allowbreak \) gives
\[
c^4\le(1-\lambda_0d)^4
\le1-(4\lambda_0-2\lambda_0^2)d
\le1-\frac75d.
\]
The cubic-plus-quartic remainder discarded here is
-\(\lambda_0^3d^3(4-\lambda_0d)\le 0,\allowbreak \) and
\(4\lambda_0-2\lambda_0^2=73839/51200>7/5.\) Therefore
\[
\begin{gathered}
N\le d(1-7d/5),\qquad\\
J_2\le\theta d(1-7d/5).
\end{gathered}\tag{R26.A.31}
\]

Equations (R26.A.28)-(R26.A.29) also give \(a\le c^2,\allowbreak \) because
2c \(\kappa>\theta/2.\) Hence (\(a-\beta)_+\le c \xi\le \theta^2.\)
As in the other rows, positivity of h and \(Eh=a-\beta\) d give \(a>0.\)

\paragraph{R26.7.3 Scalar loss and the retained second-order budget}

Let \(\lambda=1-c.\) Then \(d-13\delta_2/10\le \lambda\le d,\allowbreak \) and
\[
a\le(1-\lambda)^2-2(1-\lambda)(\delta_2-\delta_3)+\delta_3/2.
\]
This expression decreases with lambda because
\(\lambda+(\delta_2-\delta_3)\le d+\delta_2<1.\) Substitute its lower bound.
The resulting \(\delta_3\) coefficient is positive, allowing \(\delta_3/d\le \theta\) x.
The exact result is
\[
a\le1-d\left[
2-\frac{517}{320}x
-d(1-13x/10)(1-9x/80)
\right].
\]
Both factors in the last product are nonnegative. With \(d\le 1/3,\allowbreak \) the
bracket is at least
\[
\begin{gathered}
\frac53-\frac{1099}{960}x-\frac{39}{800}x^2\\
\ge\frac{2933267}{2457600}>\frac{19}{16}.
\end{gathered}
\]
Thus \(a\le 1-19d/16.\)

In (R26.A.5)-(R26.A.6), retain (R26.A.31) specifically in \(8a^2J_2\) and use the
coarser \(N\le d,\allowbreak  J_2\le \theta\) d only elsewhere. This gives
\[
\begin{gathered}
E(h^4/q)\le\\
a^4+4\theta^2a^3d\\
+a^2\left(8\theta d-\frac{51}{20}d^2\right)\\
+6\theta d^2+4ad^3.
\end{gathered}\tag{R26.A.32}
\]
The coefficient in parentheses is at least 12d/5, so substitution of
the upper bound for a is legitimate. The complete resulting polynomial is
\[
\begin{gathered}
1-\frac{215}{256}d-\frac{35269}{20480}d^2\\
+\frac{3517247}{327680}d^3-\frac{39126719}{5242880}d^4.
\end{gathered}
\]
Direct expansion gives these coefficients. The second, third and fourth
coefficients are respectively at most -\(3/2,\allowbreak \) 11,-7.
The function -\(3d/2+11d^2-7d^3\) is convex on [0, \(1/3\)], since its second
derivative is \(22-42d\ge 8.\) Its endpoint maximum is \(25/54.\) Therefore
\[
\begin{gathered}
E(h^4/q)\le1-\\
\left(\frac{215}{256}-\frac{25}{54}\right)d\\
=1-\frac{2605}{6912}d\\
\le1-\frac38d.
\end{gathered}
\]
This proves every finite \(b\ge 4.\) It would not be justified to replace
(R26.A.31) by \(N\le d\) throughout: the retained negative quadratic coefficient
in (R26.A.32) is part of the proof.

\paragraph{R26.8 The exact old C094 row and its three ordinary-fourth inputs}

The following argument proves the required bridge and its scalar estimates.

\paragraph{R26.8.1 The unweighted fourth-to-Iq bridge}

For any original positive kernel operator B, \(h=B1,\allowbreak \) and any original real
Y, row Cauchy and original-measure Holder give
\[
|BY|^2\le hB(|Y|^2),
\]
\[
I_q(BY)^3
\le E(h^4/q)\,\|B\|^4(E|Y|^4)^2.                \tag{R26.A.33}
\]
This is an ordinary L2(mu) operator norm and an ordinary fourth input.
There is no change of adjoint or input weight.

For \(B=T^2,\allowbreak  h=\bar f_2=q+r.\) Taking \(W=2r-g\) and \(z=1+2d\) in the cubic argument,
and retaining the actual row relation \(r^2\le pqd,\allowbreak \) gives
\[
E(\bar f_2^4/q)\le1-2d^2+(8-2d)\delta_2.                 \tag{R26.A.34}
\]
More explicitly,
\[
E(\bar f_2^4/q)=1-d+6\delta_2+\mathcal C+E(r^4/q),
\]
\[
\mathcal C\le d-2d^2+2(1-2d)\delta_2,\quad
E(r^4/q)\le2d\delta_2.
\]
These reproduce C094 (2.1)-(2.4).

The exact C064 \mbox{\hyperref[part:Q64.7.2]{part Q64.7.2}} norm proof is compatible with this original
space. On the cyclic subspace generated by 1 under K, split even and odd
Krylov vectors. Multiplication of \(T^2\) by the orthogonal sign operator
\(\operatorname{diag}(I_{\rm even},\allowbreak -I_{\rm odd})\) gives the self-adjoint block matrix
\[
\begin{pmatrix}
J-S_{\rm even}&1\otimes g\\
g\otimes1&S_{\rm odd}
\end{pmatrix},\qquad S=-K^2.
\]
The orthogonal complement has \(T^2=K^2\) and norm at most theta.
For \(\lambda\ge 1,\allowbreak \) the positive and negative Schur complements reduce the
only nonautomatic condition to
\[
\begin{gathered}
\lambda h_\lambda t_\lambda\le1,\quad\\
h_\lambda=\int\frac{d\nu(x)}{\lambda-x},\quad\\
t_\lambda=\int\frac{d\nu(x)}{\lambda+x}.
\end{gathered}
\]
Writing \(\alpha=int x^2/(\lambda^2-x^2)d\) nu and \(\eta=int x/(\lambda^2-x^2)d\) nu,
\[
\begin{gathered}
\lambda h_\lambda t_\lambda\\
=\frac{(1+\alpha)^2}{\lambda}-\lambda\eta^2\\
\le\frac{(1+\alpha)^2}{\lambda}.
\end{gathered}
\]
For \(k_2=(1-\theta^2)^{-1}=1024/855\) and
\(\lambda=(1+k_2\delta_2)^2,\allowbreak  \alpha\le k_2\delta_2.\) Thus
\[
\|T^2\|\le(1+1024\delta_2/855)^2.
\]
Together with (R26.A.33)-(R26.A.34), this gives exactly
\[
I_q(T^2Y)^3\le B_2(E|Y|^4)^2,
\]
\[
B_2=\{1-2d^2+(8-2d)\delta_2\}(1+1024\delta_2/855)^8.      \tag{R26.A.35}
\]

For the claimed logarithmic cost, the first factor is at most
\(R=1+(13/4)d-(45/16)d^2.\) If \(u=R-1,\allowbreak \) then
\(0\le u\le 37/48\) and \(u\ge (37/16)d.\) The elementary integral remainder inequality
\(\log(1+u)\le u-u^2/\)[2(1+u)] gives
\[
\log R\le\frac{13}{4}d-\frac{11757}{2720}d^2.
\]
For \(\lambda_2=416/855,\allowbreak  12/25<\lambda_2<1/2\) and
\[
8\log(1+\lambda_2d)
\le8\lambda_2d-\frac34d^2.
\]
Since \(11757/2720+3/4>5,\allowbreak \) the precise result is
\[
B_2\le\exp\left(\frac{24427}{3420}d-5d^2\right). \tag{R26.A.36}
\]

\paragraph{R26.8.2 Which actual messages enter this bridge}

The old row uses, for all \(b\ge 3,\allowbreak  L\ge 2,\allowbreak \)
\[
Y=pT^L(qf_b).
\]
C058 fixes this same actual direction: terminal unit arm
toward the terminal, long arm outward, inner connector toward the
internal root, internal-root unit arm outward. The two outer edges
remain directed from that root to the receiver.

The actual square-envelope estimates of Q5-C058 supply the following same-input data to Q5-C094:


\begingroup\small
\setlength{\tabcolsep}{3pt}
\setlength{\LTleft}{0pt}\setlength{\LTright}{\fill}
\begin{longtable}{@{}p{0.1200\linewidth}p{0.8000\linewidth}@{}}
\caption{All terminal pairs with minimum one at a unit root and a two-edge outlet: parameter values.}\label{tab:parameters:r2-c026:2}\\
\toprule
Entry & Actual upper envelope \(Y\le (c+N/2)^2\) \\
\midrule\endfirsthead
\toprule
Entry & Actual upper envelope \(Y\le (c+N/2)^2\) \\
\midrule\endhead
\bottomrule\endfoot
\(L\ge 4\) & \(c=1-3d/4,\allowbreak  N=R_L\); Pointwise and moment data: \(N^2\le d,\allowbreak  EN\le -d/2,\allowbreak  EN^2\le \theta d\). \\
\addlinespace[2pt]
\(L=3\) & \(c=1-4d/5,\allowbreak  N=V_3-\gamma r-q\); Pointwise and moment data: \(N^2\le (71/63)d,\allowbreak  EN^2\le \theta(71/63)d,\allowbreak  EN\le -(4/5-d)d\). \\
\addlinespace[2pt]
\(L=2\) & \(c=1-37d/64,\allowbreak  N=mR_2/\sqrt{B},\allowbreak  R_2=r+kK^2R,\allowbreak  B=1-\gamma\); Pointwise and moment data: \(N^2\le d,\allowbreak  EN\le -41d/96,\allowbreak  EN^2\le 3d/10\). \\
\end{longtable}
\endgroup


The scalar envelope base is nonnegative because it arises from a sum of
actual nonnegative factors. The L3 mean improvement uses the same R:
\(V\ge t_0^2,\allowbreak  u=kt_0\le d/(1-d/2),\allowbreak \) and
\[
\gamma\le u-u^2/2
\le d(1-d)/(1-d/2)^2\le d.
\]
For $L=2$, the original ratio $m^2/B\le9/8$ and $ER_2^2\le4d/15$ are retained, giving $EN^2\le3d/10$. These are the exact finite-envelope premises of Q5-C058 for the stated direction.

\paragraph{R26.8.3 Three compatible square moments}

Use the three actual envelopes and their same-input moment data from
paragraph R26.8.2. They are rows I, II and III of
Table~\ref{tab:unit-square-data}, for $L\ge4$, $L=3$ and $L=2$,
respectively. Their bases satisfy $c>U/2$; the compatibility
$\alpha^2\le hU^2$ also follows from the actual first and second
moments. Lemma~\ref{lem:unit-square-hermite} and the three continuous
certificates in Table~\ref{tab:unit-square-certificates} give
\begin{align}
 \E Y^4&\le(1-15d/32)^8,
 &&L\ge4,\tag{R26.A.39}\label{eq:unit-R26:fourth-long}\\
 \E Y^4&\le(1-9d/20)^8,
 &&L=3,\tag{R26.A.40}\label{eq:unit-R26:fourth-three}\\
 \E Y^4&\le(1-15d/32+3d^2/8)^8,
 &&L=2.\tag{R26.A.41}\label{eq:unit-R26:fourth-two}
\end{align}
These bounds hold on $0\le d\le1/3$ for the original direction
$Y=pT^L(qf_b)$, $b\ge3$. At $d=0$, its actual leaf integration
gives equality.

\paragraph{R26.8.7 Final exponents and the literal old union}

Use only (R26.A.39)-(R26.A.41) in (R26.A.35)-(R26.A.36). For \(L=2,\allowbreak \)
\[
\log I_q(T^2Y)^3
\le-\frac{1223}{3420}d+d^2
\le-\frac{83}{3420}d.
\]
For \(L=3\) the linear coefficient is -\(197/3420,\allowbreak \) and for \(L\ge 4\) it is
-\(1223/3420\); both also retain the nonpositive -\(5d^2\) term.
Thus the old row uniformly supplies
\[
I_q(T^2[pT^L(qf_b)])^3\le e^{-83d/3420}
\qquad(b\ge3,L\ge2).
\]
The transcription of \(B_2,\allowbreak \) all three ordinary-fourth inputs,
the inner direction and the final exponent is exact. No inverse-fourth
conclusion for this old direction follows from this argument, and none
is included in this derivation.

\paragraph{R26.9 Complete local directions, original pairs and reversal}

The complete direction table is:


\begingroup\small
\setlength{\tabcolsep}{3pt}
\setlength{\LTleft}{0pt}\setlength{\LTright}{\fill}
\begin{longtable}{@{}p{0.2200\linewidth}p{0.7000\linewidth}@{}}
\caption{All terminal pairs with minimum one at a unit root and a two-edge outlet: parameter values.}\label{tab:parameters:r2-c026:5}\\
\toprule
Entry & Inner L \\
\midrule\endfirsthead
\toprule
Entry & Inner L \\
\midrule\endhead
\bottomrule\endfoot
(1, 2) & 1; Actual bare root H: \(qA(p\bar f_2)\); New inverse fourth: \(\le 1-d/5\); \(Iq(U)^3\) bound, \(U=T^2H\): \(\le (1-d/5)^2\). \\
\addlinespace[2pt]
(1, 2) & 2; Actual bare root H: \(qA^2(p\bar f_2)\); New inverse fourth: \(\le 1-d/24\); \(Iq(U)^3\) bound, \(U=T^2H\): \(\le (1-d/24)^2\). \\
\addlinespace[2pt]
(1, 2) & 3; Actual bare root H: \(qA^3(p\bar f_2)\); New inverse fourth: \(\le 1-7d/10\); \(Iq(U)^3\) bound, \(U=T^2H\): \(\le (1-7d/10)^2\). \\
\addlinespace[2pt]
(1, 2) & 4; Actual bare root H: \(qA^4(p\bar f_2)\); New inverse fourth: \(\le 1-d/4\); \(Iq(U)^3\) bound, \(U=T^2H\): \(\le (1-d/4)^2\). \\
\addlinespace[2pt]
(1, 2) & every \(L\ge 5\); Actual bare root H: \(qA^L(p\bar f_2)\); New inverse fourth: \(\le 1-d/2\); \(Iq(U)^3\) bound, \(U=T^2H\): \(\le (1-d/2)^2\). \\
\addlinespace[2pt]
(1, 3) & 1; Actual bare root H: \(qA(p\bar f_3)\); New inverse fourth: \(\le 1-7d/10\); \(Iq(U)^3\) bound, \(U=T^2H\): \(\le (1-7d/10)^2\). \\
\addlinespace[2pt]
(1, b), every \(b\ge 4\) & 1; Actual bare root H: \(qA(p\bar f_b)\); New inverse fourth: \(\le 1-3d/8\); \(Iq(U)^3\) bound, \(U=T^2H\): \(\le (1-3d/8)^2\). \\
\addlinespace[2pt]
(1, b), every \(b\ge 3\) & every \(L\ge 2\); Actual bare root H: \(pT^L(qf_b)\); New inverse fourth: not bounded; \(Iq(U)^3\) bound, \(U=T^2H\): \(\le \exp(-83d/3420)\). \\
\end{longtable}
\endgroup


The integer partition is exhaustive: \(b=2\) uses the first five rows; if
\(b\ge 3,\allowbreak \) either \(L=1\) uses one of the next two rows or \(L\ge 2\) uses the last row.
All tests in this direction selection depend only on original integer
lengths. The analytic d intervals estimate this fixed direction.

For the new rows the literal directions are:

\begin{itemize}
\item The terminal's original unit arm points away from the terminal.
\item Its complete original b-edge arm points toward the terminal.
\item Every one of the L inner edges points from the internal root toward the terminal.
\item The internal root's one residual unit arm points toward that root.
\item Both original outer edges point from the internal root toward the receiver.
\end{itemize}

At a root, a genuine equal pair contributes \(f_l h_l\le 1.\) Full reversal
of the pair contributes \(\bar f_l W_l\le 1,\allowbreak  W_l=A \bar f_{l-1}\). These are actual
two-arm products, not arbitrary functions chosen to satisfy a budget.
Thus the new decorated message is exactly
\[
H_D=E_r^-qA^L(E_t^-p\bar f_b),\qquad U_D=T^2H_D.
\]
The terminal factor stays inside \(A^L\); the root factor stays at that root.
Original nonnegativity gives \(0<U_D\le U_*\) and therefore preserves every
new inverse-fourth bound.
The old row has, instead,
\[
H_{o,D}=E_r^+pT^L(E_t^+qf_b),\qquad U_{o,D}=T^2H_{o,D},
\]
with the exact same positivity comparison for its ordinary-fourth and
Iq bounds.

The outer operator is the actual original two-edge kernel:
\[
T^2H(x)=\sum_{z,y}\mu(z)\mu(y)
(1-K(x,z))(1-K(z,y))H(y).
\]
By antisymmetry these factors are exactly the directed edges \(y->z->x.\)
No independent row choices or tilted measures have been introduced.

Complete component reversal substitutes -K simultaneously in every
factor. It swaps p, q and A, T, leaves d unchanged, reverses every
decoration, and changes the outer \(T^2\) to \(A^2.\) Thus the new reversed
message is
\[
A^2[E_r^+pT^L(E_t^+qf_b)]
\]
with denominator p and the corresponding \(I_p\) budget. The reversed old
row is \(A^2\)[\(E_r^-qA^L(E_t^-p\bar f_b)\)] and has the corresponding old \(I_p\)
statement only. Reversing just the inner component does not effect this
global weight swap.
